This section covers running tasks on high-performance computing (HPC) clusters.
The template supports two clusters out of the box: Duke's DCC and TACC's Lonestar6,
plus local execution on your laptop/desktop.

\subsection{Environment auto-detection}

The \texttt{generic.make} file automatically detects the execution environment:

\begin{lstlisting}[language=make]
IS_TACC = 1  # when TACC_SYSTEM env var is set
IS_DCC  = 1  # when SLURM_CLUSTER_NAME = "dcc"
IS_HPC  = 1  # when either cluster is detected
\end{lstlisting}

You can override detection: \texttt{make IS\_TACC=0} forces local execution on TACC.

\subsection{SLURM basics}

SLURM is the job scheduler used by both DCC and TACC.
Key commands:
\begin{itemize}
	\item \texttt{sbatch job.slurm}: submit a batch job (returns immediately).
	\item \texttt{srun --pty bash}: start an interactive session.
	\item \texttt{squeue -u \$USER}: list your running/pending jobs.
	\item \texttt{sacct -u \$USER -S now-2days}: job history (past 2 days).
	\item \texttt{scancel <job\_id>}: cancel a job.
\end{itemize}

A typical SLURM script:
\begin{lstlisting}[language=bash]
#!/bin/bash
#SBATCH --job-name=task_name
#SBATCH --partition=common        # DCC: common, scavenger; TACC: gpu-a100
#SBATCH --time=04:00:00
#SBATCH --mem=100G
#SBATCH --cpus-per-task=4
#SBATCH --output=slurm_logs/%j.out
#SBATCH --error=slurm_logs/%j.err

# DCC: use container
singularity exec --pwd "$PWD" \
  --bind "$REPO_ROOT:$REPO_ROOT" "$IMG" \
  Rscript code/my_script.R --arg1 val1

# TACC: use conda
source /work/<project>/<user>/ls6/miniconda3/etc/profile.d/conda.sh
conda activate my_env
python code/my_script.py --arg1 val1
\end{lstlisting}

\subsection{Container execution (DCC)}

On DCC, R tasks run inside a Singularity container to ensure reproducible package versions.
The container image (\texttt{.sif} file) is stored at a shared path.

Key flags for \texttt{singularity exec}:
\begin{itemize}
	\item \texttt{--pwd "\$PWD"}: set working directory inside the container.
	\item \texttt{--bind "\$REPO\_ROOT:\$REPO\_ROOT"}: mount the repository so symlinks resolve.
\end{itemize}

The \texttt{generic.make} wrappers \texttt{RUN\_CMD} and \texttt{RSCRIPT} handle this automatically.

\subsection{Conda execution (TACC)}

On TACC, conda environments manage package dependencies.
Conda lives under a project-specific path on \texttt{/work/}, e.g.:
\begin{lstlisting}[language=bash]
/work/<project_num>/<user>/ls6/miniconda3/
\end{lstlisting}
\noindent not in the home directory (which has limited quota).

\subsection{Job dependency chaining}

For multi-step pipelines within a single task, use SLURM dependencies:

\begin{lstlisting}[language=bash]
job1=$(sbatch --parsable step1.slurm)
job2=$(sbatch --parsable --dependency=afterok:$job1 step2.slurm)
job3=$(sbatch --parsable --dependency=afterok:$job2 step3.slurm)
\end{lstlisting}

The \texttt{slurm\_helpers.sh} script provides functions for this:
\texttt{submit\_slurm\_job}, \texttt{submit\_job\_chain},
and \texttt{wait\_for\_jobs}.

\subsection{Two execution patterns}

\textbf{Pattern 1: Submit \& Continue.}
The Makefile submits SLURM jobs and returns immediately.
Downstream jobs use \texttt{--dependency=afterok} to chain.
The user re-runs \texttt{make} after jobs complete.
This is the most common pattern.

\textbf{Pattern 2: Submit \& Wait.}
The Makefile submits a job and blocks until completion (using \texttt{sbatch --wait}
or polling via \texttt{squeue}).
After completion, local post-processing runs.
Use this when the entire pipeline must complete in one \texttt{make} invocation.

\subsection{Data synchronization}

The \texttt{sync\_from\_cluster.sh} script downloads task outputs from HPC to your local machine:

\begin{lstlisting}[language=bash]
./sync_from_cluster.sh --tacc task_name     # Download output/ from TACC
./sync_from_cluster.sh --dcc task_name      # Download output/ from DCC
./sync_from_cluster.sh --list --tacc        # List available outputs
./sync_from_cluster.sh --dry-run task_name  # Preview transfer
\end{lstlisting}

After syncing, run \texttt{make} for downstream tasks to recreate local symlinks.

\subsection{Interactive sessions}

For debugging or iterating on code:
\begin{lstlisting}[language=bash]
# DCC: interactive container session
make -C tasks/<task_name> INTERACTIVE=1

# TACC: interactive node
srun -p gpu-a100 --time=02:00:00 --pty bash
source /work/<project>/<user>/ls6/miniconda3/etc/profile.d/conda.sh
conda activate my_env
\end{lstlisting}
